% Pursuing Stacks (cotes 134-2 to 134-8): archive cote 134-2, site page 40 (PDF page 41).
% appendix (letters to Breen). Draft compared: ps1.tex:1195-1203.
% AI-transcribed by Claude Opus 5.5 (claude-opus-5-5, adaptive thinking, effort high), Message Batches API, batch ps1 (msgbatch_01Rv5hCaNTtUt4JHx1yqoHPL), 2026-11-12; not yet checked by eye.
\scanpage{134-2/40}
\mspage{\hand{25}}

dont les invariants sont $H^3(B_G \text{ mod } X,N)$, $H^2(B_G \text{ mod } X,N)$, et $Z^1(G,N)$ (groupe
des 1-cocyles de $G$ à coéfficients dans $N$)- je te laisse le soin de deviner
quel est le complexe tronqué qui le décrit ! J'ai écrit il y a quelques
mois à Deligne à ce sujet, et l'ai prié de t'envoyer une copie de \xout{\unsure{envoyer}}
la lettre.

4) Soit $X$ un topos localement annelé, on peut considérer les \uline{Algèbres
d'Azumaya sur} $X$ (i.e. les Algèbres localement isomorphes à une alg. de
matrices d'ordre $n$, $n \geq 1$) comme les objets d'une 2-catégorie de Picard,
où la catégorie $\underline{\mathrm{Hom}}(A,B)$, pour $A$ et $B$ des Algèbres d'Azumaya, est la caté-%
gorie des ``trivialisations'' de $A^o\otimes B$, i.e. des couples $(E,\emptyset)$, $E$ un Module
localement libre et $\emptyset$ un isomorphisme $\underline{\mathrm{End}}(E)\simeq A^o\otimes B$. Il faut travailler
un peu pour définir les accouplements\xout{x} $\underline{\mathrm{Hom}}(A,B)\mathrm{x}\underline{\mathrm{Hom}}(B,C)\hand{\to}\underline{\mathrm{Hom}}(A,C)$; l'opé-%
ration $\otimes$ dans la 2-catégorie de Picard à construire est bien sûr le produit
tensoriel d'Algèbres, et l'opération ``puissance $-1$'' est le passage à l'al-%
gèbre opposée. On \unsure{vérfie|vérifie}\xout{,xx passant aux champs correspondant, qui sont} qu'en
associant à toute Algèbre d'Azumaya la 1-gerbe de ses trivialisations, on
trouve un 2-$\otimes$-foncteur de la 2-catégorie de Picard (dite ``de Brauer'') dans
celle des 1-gerbes liées par $\underline{G}_m$ , qui est 2-fidèle. Les invariants de
la première sont donc les groupes $H^2(X,\underline{G}_m)_{Br}$ , $H^1(X,\underline{\unsure{H|G}}_m)$ et $H^o(X,\underline{G}_m)$ , où
dans le premier terme l'indice Br désigne le sous-groupe du $H^2$ formé des
classes de cohomologie provenant d'Algèbres d'Azumaya. On aurait envie
de parler du 2-champ de Picard des Algèbres d'Azumaya sur des objets variab
les de $X$, mais c'est bien une 2-catégorie de Picard fibrée sur $X$, mais pas
tout à fait un 2-champ (whatever that means), sans doute - la condition de
2-recollement (whatever that means) ne doit pas être satisfaite - sinon il
n'y aurait pas d'indice Br au $H^2$ ...
\margin{espace}

\inserted{(App.2)} La considération des $n$-catégories de Picard \inserted[typed]{strictes} (qui s'imposent à nous pas
à pas dans un contexte essentiellement ``commutatif'') me semblent la clef du
passage de l'algèbre homologique ordinaire (``commutative''), en termes de
complexes, à une algèbre homologique non commutative, du fait qu'elles
donnent une interprétation géométrique correcte des ``complexes tronqués à
l'ordre $n$'' \inserted{(en tant \unsure{qu'objets|qui objets} de catégories dérivées)}, donc, essentiellement (par passage à la limite sur $n$) des
complexes tout courts. L'idée naive qui se présente est alors que les
``complexes non commutatifs'' (qui seraient les objets-fantômes d'une algèbre
homologique non commutative) sont peut-être ce qui reste des $n$-catégories
de Picard (strictes) quand on oublie \unsure{leut|leur} caractère additif, i.e. leur
structure de Picard - c'est à dire qu'on ne retient que la $n$-catégorie !
\inserted{(Quand on se place sur un topos $X$, on s'intéresse donc aux $n$-champs sur $X$...)}
A vrai dire, cette idée est venue d'abord d'une autre direction, quand il
s'est agi en géométrie algébrique de démontrer des \uline{\deleted{formules}\inserted{théorèmes} de Lefschetz}

\margin{\hand{40}}
